

%  Información Complementaria para la comprensión, que estaría de más en el Informe  %


    \section{Instrumental Utilizado} \label{ap:Instrumental}
Para las prácticas experimentales realizadas en este trabajo se contaba con el siguiente equipo:
        \subsection{Dispositivo Experimental}
            \subsubsection{Osciloscopio}
Un osciloscopio Tektronix TBS1104 digital de cuatro canales ($100 \n{MHz}$, $1 \n{GS}/\tx{s}$).
            \subsubsection{Fuente De Corriente Alterna}
Una fuente de corriente alterna PROTOMAX HY3005D-3 de alta tensión ($220 \n{V}$, $50 \n{Hz}$).
            \subsubsection{Fuente De Corriente Continua}
Un generador de funciones Sinometer VC2002 ($220 \n{V}$, $50 \n{Hz}$).
            \subsubsection{Amplificador}
Un amplificador de potencia profesional SKP MAX G1400 ($1400 \n{W}$).
            \subsubsection{Bomba Peristáltica}
Una bomba peristáltica APEMA PC25-10-F-S de Clase 1 ($10 \n{W}$).
            \subsubsection{Agitador Magnético}
Un agitador Magnético ArcanO HJ-3 ($2 \n{A}$).
            \subsubsection{Espectrofotómetro}
Un espectrofotómetro BioTraza Modelo 721 ($40 \n{W}$).
        \subsection{Muestra a tratar}
            \subsubsection{Azul de Metileno}
Azul de metileno ($\tx{C}_{16}\tx{H}_{18}\tx{ClN}_3\tx{S}$) del laboratorio Cicarelli en $25 \n{g}$.
    \section{Resultados Obtenidos} \label{ap:Resultados}
